% Lean: EconLean/Macro/Climate/SocialCostCarbon.lean
\chapter{The Social Cost of Carbon}
\label{chap:climate}

Sources: Nordhaus (DICE); Golosov, Hassler, Krusell and Tsyvinski (2014). Log utility, constant saving rate $s$,
damages $e^{-\gamma S}$ and exogenous damage-free output $\tilde Y$.

\begin{definition}[Carbon cycle and welfare]
\label{def:scc_model}
\lean{Econ.Climate.Params, Econ.Climate.Params.remain, Econ.Climate.Params.outputAfter, Econ.Climate.Params.welfareChange, Econ.Climate.Params.kernel, Econ.Climate.Params.scc}
\leanok
A unit emitted at date 0 leaves $r_j = \phi_L + (1-\phi_L)\phi_0(1-\phi)^j$ in the atmosphere after $j$ periods. The
welfare change from $e$ extra units is $\Delta W(e) = \sum_j\beta^j(\log C_j(e) - \log C_j(0))$. The SCC is
$-C_0\,\Delta W'(0)$.
\end{definition}

\begin{theorem}[GHKT formula]
\label{thm:scc}
\lean{Econ.Climate.Params.summable_kernel, Econ.Climate.Params.tsum_kernel, Econ.Climate.Params.welfareChange_eq, Econ.Climate.Params.hasDerivAt_welfareChange, Econ.Climate.Params.scc_eq}
\leanok
\uses{def:scc_model}
$\sum_j\beta^jr_j = \phi_L/(1-\beta) + (1-\phi_L)\phi_0/(1-\beta(1-\phi))$, $\Delta W(e) = -\gamma e\sum_j\beta^jr_j$, and
$\mathrm{SCC} = \gamma C_0\sum_j\beta^jr_j$.
\end{theorem}
\begin{proof}
\leanok
Each log difference equals $-\gamma r_je$ because $\tilde Y$ and $s$ cancel. Sum the two geometric series.
\end{proof}

\begin{theorem}[GHKT with capital]
\label{thm:scc_capital}
\lean{Econ.Climate.Params.logYcap, Econ.Climate.Params.welfareCap, Econ.Climate.Params.sccCap, Econ.Climate.Params.logYcap_sub, Econ.Climate.Params.tsum_cauchy, Econ.Climate.Params.welfareCap_eq, Econ.Climate.Params.sccCap_eq}
\leanok
\uses{thm:scc}
With $Y_j = e^{-\gamma S_j}A_jK_j^\alpha$, full depreciation and saving rate $\alpha\beta$, an extra unit of
emissions lowers $\log Y_j$ by $\gamma\sum_{k\le j}\alpha^{j-k}r_k$, and $\mathrm{SCC} = \gamma Y_0\sum_j\beta^jr_j$.
\end{theorem}
\begin{proof}
\leanok
Induct on $j$ using $y_{j+1} = -\gamma(S_{j+1}+r_{j+1}e) + a_{j+1} + \alpha(\log\alpha\beta + y_j)$. The discounted
sum is a Cauchy product of $\sum(\alpha\beta)^m$ and $\sum\beta^kr_k$, which equals $\mathrm{kernel}/(1-\alpha\beta)$.
Multiplying by $C_0 = (1-\alpha\beta)Y_0$ cancels the factor.
\end{proof}

\begin{theorem}[Discounting]
\label{thm:scc_discount}
\lean{Econ.Climate.Params.kernel_pos_of, Econ.Climate.Params.kernel_mono, Econ.Climate.Params.kernel_unbounded}
\leanok
\uses{thm:scc}
The kernel is increasing in $\beta$. If $\phi_L > 0$ it is positive and unbounded as $\beta \to 1$.
\end{theorem}
\begin{proof}
\leanok
Both denominators fall as $\beta$ rises. Taking $\beta = 1-\varepsilon$ gives a kernel of at least $\phi_L/\varepsilon$.
\end{proof}
